% Introduction --- derived from the Avance7 "Problem Statement" and EDA
% findings, reorganized as a paper introduction with an explicit thesis
% and contributions.
\section{Introduction}
\label{sec:intro}

Efficient agricultural management requires precise knowledge of what is
grown, where, and in what quantity. Traditional methods based on
physical field inspection are expensive, slow, and difficult to scale
across large territories. Sentinel-2 imagery provides continuous,
freely available coverage, yet turning it into reliable crop maps is a
hard machine-learning problem for three structural reasons.

\paragraph{Severe class imbalance.}
Of the 18 crop classes in PASTIS-R~\cite{garnot2021pastis}, meadow
accounts for up to 30\% of all parcels, while agronomically valuable
classes such as rapeseed, sorghum, and orchard are minority classes
(below 1\%). A model that ignores rare classes has little commercial
value, so macro-averaged metrics are the right objective.

\paragraph{Intra-class spectral ambiguity.}
Crops such as soft winter wheat and winter barley are virtually
indistinguishable in a single RGB image. Discrimination requires
exploiting the \emph{temporal dynamics} of the spectral signature
across the year --- greenup onset, NDVI peak, harvest date --- which
motivates time-series architectures over single-image
classifiers~\cite{tarasiou2023tsvit,garnot2021utae}.

\paragraph{Production constraints.}
A deployable system must produce parcel-level predictions (not only
pixel-level), be explainable to non-technical users, and operate with
inference latencies compatible with an interactive interface.

\paragraph{Thesis.}
We argue that crop-area quantification is best served by a
\emph{perceiver/reasoner} architecture. A perception stack of
specialised models --- temporal transformers, tabular classifiers over
AlphaEarth annual embeddings~\cite{brown2025alphaearth}, and a
contrastive vision-language branch~\cite{li2025farslip} --- produces
auditable, plain-text observations per parcel. A \emph{frozen} frontier
LLM then reasons over those observations and communicates with the
user, following the \emph{Be My Eyes}
pattern~\cite{huang2025bemyeyes}. The language layer is a
communication, explanation, and tool-orchestration component --- not a
pixel classifier --- which keeps every cited figure traceable to a
deterministic tool result and preserves the option of on-premises
operation for data sovereignty.

\paragraph{Contributions.}
\begin{enumerate}[leftmargin=*]
    \item A comparative study of dense segmentation architectures on
    PASTIS-R, identifying TSViT-pheno~\cite{tarasiou2023tsvit} as the
    best individual model (mIoU 0.6253, macro F1 0.7500 on the
    held-out fold), and heterogeneous ensembles that improve the
    out-of-fold macro F1 over any single member
    (Section~\ref{sec:results}).
    \item A grounded conversational layer built on the \emph{Be My
    Eyes} pattern with a dual frozen reasoner (Gemini~3.5~Flash cloud /
    Qwen3-30B-A3B on-premises) and a closed geospatial toolset with
    Spatial-RAG (Section~\ref{sec:method}).
    \item An honest study of spatial transferability from France to
    Catalonia (Sen4AgriNet~\cite{sykas2022sen4agrinet}) and a few-shot
    $k$-shot curve on EuroCropsML~\cite{reuss2025eurocropsml},
    showing that zero-shot transfer fails and that few-shot adaptation
    is required, in line with reported
    limits~\cite{harvesting2026alphaearth}
    (Section~\ref{sec:results}).
\end{enumerate}

The working dataset is PASTIS-R (metropolitan France), with 2{,}433
multi-temporal Sentinel-2 patches of $128\times128$ pixels, real
farmer-validated labels, 18 crops, and spatially disjoint folds that
prevent information leakage between training and evaluation.
